\documentclass[11pt]{article}
\usepackage[margin=1in]{geometry}
\usepackage[T1]{fontenc}
\usepackage{lmodern}
\usepackage{amsmath,amssymb,amsthm,mathrsfs}
\usepackage{graphicx}
\usepackage{microtype}
\usepackage[colorlinks=true,linkcolor=blue,citecolor=blue,urlcolor=blue]{hyperref}

\newtheorem{theorem}{Theorem}
\newtheorem{corollary}[theorem]{Corollary}
\newtheorem{lemma}[theorem]{Lemma}
\newcommand{\R}{\mathbb R}
\newcommand{\Schat}{\mathcal S}
\newcommand{\FW}{\mathcal F_W}
\newcommand{\Fs}{\mathcal F_\sigma}
\newcommand{\Tm}{\mathfrak T}
\newcommand{\Mm}{\mathfrak M}

\title{Quadratic chirps are not trace-class\\Fourier--Wigner multipliers}
\author{Candidate partial result for arXiv:2502.16248}
\date{11 August 2026}

\begin{document}
\maketitle

\begin{center}
\fbox{\parbox{0.91\linewidth}{\textbf{Status.} Candidate substantial partial
result, likely valid, requiring human review.  The source's explicit
nonmeasure candidate is ruled out quantitatively.  The full classification of
trace-class Fourier--Wigner multipliers remains open.}}
\end{center}

\section{Source question and result}

For $z=(x,\xi)\in\R^{2d}$, let $\rho(z)$ be the symmetric time--frequency
shift and let
\[
  \FW(T)(z)=\operatorname{tr}(T\rho(-z))
\]
be the Fourier--Wigner transform.  For $m\in L^\infty(\R^{2d})$, the associated
Fourier--Wigner multiplier is formally
\[
  \Tm_m(T)=\FW^{-1}\bigl(m\FW(T)\bigr).
\]
The algebra $\Mm_1(\R^{2d})$ consists of those $m$ for which this map is
bounded on the trace class $\Schat^1$.

Samuelsen proves that $\Fs(\mathscr M(\R^{2d}))\subseteq\Mm_1$ and asks
whether equality holds \cite[Question~1]{Samuelsen}.  The paper singles out
\[
                  s(z)=\sin(\pi|z|^2)
\]
as a possible counterexample: $\Fs(s)$ is not a finite measure, but the source
leaves membership of $s$ in $\Mm_1$ undecided.  The theorem below rules out
this candidate and, more generally, every nonconstant radial quadratic chirp.

\begin{center}
\includegraphics[width=0.91\linewidth]{figures/open_question_crop.png}

\small\textit{Figure 1. The proposed sine candidate and Question 1 in
\cite[p.~21]{Samuelsen}, cropped from the local source PDF.}
\end{center}

\section{Idea of the proof}

Test the multiplier on rank-one projections generated by increasingly
squeezed Gaussians.  Their Fourier--Wigner transforms are anisotropic
Gaussians of fixed $L^2$ mass.  Multiplication by a unimodular quadratic chirp
therefore leaves the Hilbert--Schmidt norm of the output equal to one.

On the other hand, the output integral kernel can be evaluated exactly.  It is
narrow in the difference variable $t-s$, with width of order $r^{-1/2}$.
Schur's test makes its operator norm decay like $r^{-d/2}$.  Since
\[
              \|A\|_2^2\le \|A\|_\infty\|A\|_1,
\]
fixed Hilbert--Schmidt mass and vanishing operator norm force the trace norm to
grow like $r^{d/2}$.  For the sine, the two conjugate chirps cannot cancel:
an exact oscillatory Gaussian integral leaves a uniformly positive amount of
Hilbert--Schmidt mass.

\section{The squeezed-Gaussian calculation}

For $r>0$, define the normalized Gaussian
\[
  g_r(t)=(2r)^{d/4}e^{-\pi r|t|^2},\qquad t\in\R^d,
\]
and the rank-one operator $G_r=g_r\otimes g_r$.  Thus
$\|G_r\|_1=1$.  Direct calculation with the source conventions gives
\begin{equation}\label{eq:ambiguity}
 \FW(G_r)(x,\xi)
 =\exp\left[-\frac\pi2\left(r|x|^2+r^{-1}|\xi|^2\right)\right].
\end{equation}

We use the following centered-kernel formula.  If $F=\FW(A)$ and $K_A$ is
the integral kernel of $A$, then
\begin{equation}\label{eq:centered}
 K_A(t,s)=\int_{\R^d}F(t-s,\xi)
       e^{2\pi i\xi\cdot(t+s)/2}\,d\xi.
\end{equation}
This follows immediately by inserting the definition of $\rho(x,\xi)$; it is
also the usual ambiguity-function/kernel relation.

\begin{lemma}[Chirp kernel bound]\label{lem:kernel}
Let $c\in\R\setminus\{0\}$ and $q_c(z)=e^{\pi i c|z|^2}$.  Put
$A_{r,c}=\Tm_{q_c}(G_r)$, initially interpreted in $\Schat^2$.  Then
\begin{align}
 \|A_{r,c}\|_2&=1,\label{eq:hsone}\\
 \|A_{r,c}\|_\infty&\le
       \left(\frac{2}{|c|r}\right)^{d/2}.\label{eq:opbound}
\end{align}
\end{lemma}

\begin{proof}
Equation \eqref{eq:hsone} follows from the unitarity of $\FW:\Schat^2\to
L^2(\R^{2d})$, the unimodularity of $q_c$, and
$\|\FW(G_r)\|_2=\|G_r\|_2=1$.

For the operator norm, set
\[
    D_{r,c}=\frac1{2r}-ic,\qquad
    u=\frac{t+s}{2},\qquad x=t-s.
\]
Applying \eqref{eq:centered} to $q_c\FW(G_r)$ and evaluating the Gaussian
Fourier integral gives, with the square-root branch determined by
$\operatorname{Re}D_{r,c}>0$,
\begin{equation}\label{eq:kernel}
 K_{r,c}(t,s)=D_{r,c}^{-d/2}
 \exp\left(-\frac{\pi|u|^2}{D_{r,c}}
           -\frac{\pi r}{2}|x|^2+\pi ic|x|^2\right).
\end{equation}
Since $\operatorname{Re}(D_{r,c}^{-1})>0$ and
$|D_{r,c}|\ge|c|$,
\[
 |K_{r,c}(t,s)|
 \le |c|^{-d/2}e^{-\pi r|t-s|^2/2}.
\]
Therefore
\[
 \sup_t\int_{\R^d}|K_{r,c}(t,s)|\,ds,
 \quad
 \sup_s\int_{\R^d}|K_{r,c}(t,s)|\,dt
 \le |c|^{-d/2}\left(\frac2r\right)^{d/2}.
\]
Schur's test proves \eqref{eq:opbound}.
\end{proof}

\begin{theorem}[Quadratic chirp obstruction]\label{thm:chirp}
For every $c\in\R\setminus\{0\}$,
\[
        q_c(z)=e^{\pi ic|z|^2}\notin\Mm_1(\R^{2d}).
\]
More precisely, whenever $A_{r,c}$ is trace class,
\begin{equation}\label{eq:chirplower}
 \|A_{r,c}\|_1\ge
       \left(\frac{|c|r}{2}\right)^{d/2}.
\end{equation}
\end{theorem}

\begin{proof}
For any trace-class operator $A$, its singular values give
$\|A\|_2^2\le\|A\|_\infty\|A\|_1$.  Lemma~\ref{lem:kernel} therefore yields
\eqref{eq:chirplower}.  Since $\|G_r\|_1=1$ and the right-hand side tends to
infinity, $\Tm_{q_c}$ cannot be bounded on $\Schat^1$.
\end{proof}

\section{The source's sine candidate}

\begin{theorem}[Negative decision for the proposed candidate]\label{thm:sine}
Let $s(z)=\sin(\pi|z|^2)$.  Then
\[
                   \boxed{s\notin\Mm_1(\R^{2d}).}
\]
Quantitatively, if $B_r=\Tm_s(G_r)$ and $B_r$ is trace class, then
\begin{equation}\label{eq:sinelower}
 \|B_r\|_1\ge
 \frac{1-5^{-d/2}}2\left(\frac r2\right)^{d/2}.
\end{equation}
\end{theorem}

\begin{proof}
Because $s=(q_1-q_{-1})/(2i)$, Lemma~\ref{lem:kernel} gives
\begin{equation}\label{eq:sineop}
                    \|B_r\|_\infty\le(2/r)^{d/2}.
\end{equation}
By Plancherel and \eqref{eq:ambiguity},
\[
 \|B_r\|_2^2
 =\int_{\R^{2d}}\sin^2\!\bigl(\pi(|x|^2+|\xi|^2)\bigr)
   e^{-\pi(r|x|^2+r^{-1}|\xi|^2)}\,dx\,d\xi.
\]
Using $\sin^2\theta=(1-\operatorname{Re}e^{2i\theta})/2$ and the complex
Gaussian integral, this equals
\begin{equation}\label{eq:exacths}
 \frac12\left[1-\operatorname{Re}\bigl(
 (r-2i)^{-d/2}(r^{-1}-2i)^{-d/2}\bigr)\right].
\end{equation}
The oscillatory factor in braces has modulus
\[
 \bigl((r^2+4)(r^{-2}+4)\bigr)^{-d/4}
 =\bigl(17+4(r^2+r^{-2})\bigr)^{-d/4}
 \le5^{-d/2}.
\]
Hence $\|B_r\|_2^2\ge(1-5^{-d/2})/2$.  Combining this with
\eqref{eq:sineop} and
$\|B_r\|_2^2\le\|B_r\|_\infty\|B_r\|_1$ proves
\eqref{eq:sinelower}, which diverges with $r$.
\end{proof}

\begin{corollary}\label{cor:finite}
No nonconstant finite linear combination of distinct radial quadratic chirps
$z\mapsto e^{\pi ic_j|z|^2}$ belongs to $\Mm_1(\R^{2d})$.
\end{corollary}

\begin{proof}
Subtract the constant term, which is already a measure multiplier.  For the
remaining distinct nonzero $c_j$, the triangle inequality in
\eqref{eq:opbound} makes the output operator norm $O(r^{-d/2})$.  Expanding
the squared Hilbert--Schmidt norm produces diagonal terms $\sum_j|a_j|^2$;
each cross term is a complex Gaussian integral with nonzero chirp parameter
and tends to zero as $r\to\infty$.  Thus the Hilbert--Schmidt norm stays
bounded below, and the same singular-value argument forces trace-norm
divergence.
\end{proof}

\section{Scope, upgrade attempts, and search status}

Theorem~\ref{thm:sine} settles the concrete candidate immediately preceding
the source's Question~1, but it does not answer the universal classification.
The strongest surviving reformulation is operator-algebraic: a completely
bounded Weyl-covariant multiplier has a Fourier--Stieltjes symbol, so a full
affirmative answer would follow if every bounded multiplier on $\Schat^1$
were automatically completely bounded.  Covariance alone does not make this
formal; finite two-torsion Weyl systems have diagonal transposition maps with
bounded norm one and growing completely bounded norm.

Two further upgrade routes were tested.  Compactly supported localization
reduces to the classical $L^1$ multiplier theorem by the source's main result,
but its comparison constants grow under exhaustion.  A Lagrangian-restriction
approach likewise runs into the lack of uniform equicontinuity of the
Fourier--Wigner images of the trace-class unit ball.  The squeezed-Gaussian
calculation identifies the same concentration obstruction quantitatively.
The accompanying attempt record gives the detailed route audit.

\paragraph{Bounded novelty check.}
The run indexes, exact question, exact sine candidate, quadratic-chirp terms,
Weyl-covariant-map terminology, twisted Fourier multipliers, and citing works
were searched through 11 August 2026.  No later full answer and no prior
statement of Theorems~\ref{thm:chirp}--\ref{thm:sine} was found.  The 1984
framework of Werner \cite{Werner} characterizes positive covariant
correspondence rules, not arbitrary bounded multipliers.  This is a bounded
novelty audit, not a definitive priority claim.

\paragraph{Human-review recommendation.}
First check the centered kernel convention in \eqref{eq:centered}.  Then audit
the factor $D_{r,c}^{-d/2}$ and the two uses of
$\|A\|_2^2\le\|A\|_\infty\|A\|_1$.  Finally verify the modulus estimate after
\eqref{eq:exacths}.  These are the only normalization-sensitive points.

\begin{thebibliography}{9}
\bibitem{Samuelsen}
H.~J. Samuelsen,
\emph{Fourier--Wigner multipliers and the Bochner--Riesz conjecture for
Schatten class operators}, arXiv:2502.16248v2 (2025).

\bibitem{Werner}
R.~F. Werner,
\emph{Quantum harmonic analysis on phase space},
J. Math. Phys. \textbf{25} (1984), 1404--1411.
\end{thebibliography}

\end{document}
